\documentclass[11pt]{article}

\usepackage[margin=1in]{geometry}
\usepackage{amsmath}
\usepackage{amssymb}
\usepackage[T1]{fontenc}
\usepackage{textcomp}
\usepackage{booktabs}
\usepackage{hyperref}
\usepackage{parskip}

\title{Spring-constant scales and units in elastic network models}
\author{Notes for penm}
\date{2026-10-15}

\begin{document}

\maketitle

\section{The question}

In a harmonic network with $k_{ij} = \gamma\,\kappa_{ij}$, where $\gamma$ is an overall
scale and $\kappa_{ij}$ the model's shape, temperature and stiffness enter only through
$\beta\gamma$, in \AA$^{-2}$.
Mode shapes, structural responses $\mathbf{K}^{+}\mathbf{f}$ and differences of
$\beta T S$ between two networks are independent of both.
Fluctuations scale as $(\beta\gamma)^{-1}$, $\beta V$ as $\beta\gamma$, and absolute entropic
terms $\ln(\beta\lambda/2\pi)$ depend on $\beta\gamma$ and on the length unit.
A free energy that adds an energy to an entropic term therefore needs $\gamma$ in physical
units, consistent with the units of $\beta$.

These notes record what five papers say about the scale $\gamma$, its units, what it was
fitted to, and at what temperature.

\section{The papers}

\subsection*{Ming \& Wall 2005 \cite{ming2005}}

Fits C$_\alpha$ networks of trypsinogen to a CHARMM22 all-atom harmonic model (vacuum
minimisation, Hessian computed in CHARMM).
Two targets: the Kullback--Leibler divergence $D$ between the C$_\alpha$ distributions,
and the C$_\alpha$ mean-square displacements (MSDs).
The optimal $\gamma$ for a given network follows in closed form from $D$ (their eq.~10).

\begin{center}
\begin{tabular}{lllll}
\toprule
model & target & $\gamma$ (kcal/mol/\AA$^2$) & $r_c$ (\AA) & $\epsilon$ \\
\midrule
ANM  & $D$   & 53.4 & 7.75 & --- \\
ANM  & MSDs  & 0.4  & 15.4 & --- \\
BENM & $D$   & 4.26 & 10.5 & 42 \\
BENM & MSDs  & 1.2  & 10.5 & 42 \\
\bottomrule
\end{tabular}
\end{center}

BENM multiplies the spring constant of backbone neighbours by $\epsilon$.
The reference is harmonic, so its fluctuations scale with $k_BT$ exactly as the network's
do: these $\gamma$ are mechanical stiffnesses, independent of temperature (the MSDs are
quoted at 300~K).
The $D$-optimal ANM ``exhibits poor MSDs''; fitting the MSDs instead moves $\gamma$ by a
factor of about 130 and the cutoff from 7.75 to 15.4~\AA.
One protein.

\subsection*{Orellana et al.\ 2010 \cite{orellana2010}}

ED-ENM, fitted to essential dynamics of MD metatrajectories (AMBER, CHARMM, OPLS/AA;
$\geq 10$~ns; six training proteins, 32+ test proteins).
Sequence neighbours up to $i+3$ get $K_{ij} = C_{\mathrm{seq}}/S_{ij}^{2}$; other pairs within
a size-dependent cutoff (about 8--16~\AA) get an inverse sixth power of distance.
The constants are $C_{\mathrm{seq}} = 60$ and $C_{\mathrm{cart}} = 6$~kcal/(mol\,\AA$^2$),
found by matching total variance, variance profile and reduced variance of the MD essential
modes, and robust to $\pm 10$ and $\pm 1$ respectively.
The long-range term is written $C_{\mathrm{cart}}/d_{ij}^{6}$ in eq.~6 and
$(C_{\mathrm{cart}}/d_{ij})^{6}$ in eq.~8; only the first is consistent with the stated
units.
The energy is $E = \sum K_{ij}(R_{ij}-R^0_{ij})^2$, without the factor $\tfrac12$.
The MD temperature is not stated in the main text.
Calibrating against MD gives effective constants at the MD temperature, through
$K^{\mathrm{app}}_{ij} = k_BT/\langle (R_{ij}-R^0_{ij})^2\rangle$ and $K_\nu = k_BT/\lambda$.

\subsection*{Leioatts et al.\ 2012 \cite{leioatts2012}}

Fits several spring functions (cutoff, exponential, HCA-like, bonded/nonbonded) to
$>1~\mu$s MD of three GPCRs.
The inverted ENM eigenvalues ``are normalized \ldots\ such that the total amplitude of
predicted motion is identical'' to MD, so the fits determine the shape of $k(r)$, not its
scale; the cutoff model uses $k = 1$.
The potential is $U_{ij} = k(R^\circ_{ij})(R_{ij}-R^\circ_{ij})^2$, again without
$\tfrac12$.
Main result: after fitting, all spring functions perform alike, because all make the first
shell of neighbours (about 3.8~\AA) much stiffer than more distant contacts, by about an
order of magnitude or more.

\subsection*{Fuglebakk et al.\ 2013 \cite{fuglebakk2013}}

Benchmarks ANM (cutoffs 0.8 and 1.8~nm), GNM, HCA, pfANM, REACH and $\beta$GM against
covariances from MD for seven proteins (five run at 300~K; two taken from earlier
publications).
Every covariance is normalised by its trace, which ``removes a global amplitude factor from
the fluctuations and covariances, which none of the ENMs describes realistically anyway'';
of $\gamma$: ``its value is not important''.
No scale is determined.
Results bearing on parameters:
\begin{itemize}
\item the ANM cutoff that best reproduces MD covariances lies in 0.8--1.0~nm; cutoffs that
best reproduce B-factors are much longer and suppress collective motion;
\item models with weak long-range interactions (HCA, REACH, ANM$_{0.8}$, $\beta$GM) beat
those with strong ones (ANM$_{1.8}$, pfANM); the precise decay law matters little;
\item B-factors are discouraged as a benchmark, and ``are not necessarily proportional to the
temperature''.
\end{itemize}
REACH is used with parameters from MD of a myoglobin dimer.

\subsection*{Dubanevics \& McLeish 2022 \cite{dubanevics2022}}

ANM and BENM of three proteins (CAP, GST, Mpro) compared with 300~K MD (PCA of
mass-weighted covariance).
Potential $V_{ij} = \tfrac{k}{2}(r_{ij}-r^{(0)}_{ij})^2$ within $r_c$, starting from
$k = 1$~kcal/(mol\,\AA$^2$), then rescaled so that the first ENM mode frequency matches
the first MD mode (their eq.~3, as printed):
\begin{equation*}
k_{\mathrm{scaled}} = k\,\frac{\omega_{\mathrm{MD}}(1)}{\omega_{\mathrm{ENM}}(1)} .
\end{equation*}
The preceding sentence states $\omega^2 \propto k$, so matching the frequencies requires the
square of this ratio; whether the tabulated values used the printed or the squared form
cannot be told from the paper.
Frequencies come from the mass-weighted Hessian and enter a quantum harmonic free energy,
so the absolute scale matters there.
MD: 300~K, 300~ns production, PCA on the last 50--150~ns (their Table~5).

\begin{center}
\small
\begin{tabular}{llll}
\toprule
protein & $r_c$ (\AA) & ANM ($\epsilon=1$) & BENM ($\epsilon=100$) \\
\midrule
CAP  & 8.0 & 31.10 & 23.56 \\
GST  & 8.5 & 25.71 & 20.55 \\
Mpro & 8.5 & 12.28 & 9.01 \\
\bottomrule
\end{tabular}
\end{center}
($k_{\mathrm{scaled}}$ in kcal/mol/\AA$^2$, their Table~3.)
The optimal $\epsilon$ is protein-dependent (100 for CAP and GST, 200 for Mpro), and the
low:high frequency ratio across the BENM gap is 2:1, as in Ming \& Wall.

\section{Observations}

\begin{enumerate}
\item \textbf{Most comparisons are scale-free by construction.}
Leioatts and Fuglebakk normalise amplitudes before comparing; they constrain the shape of
$k(r)$ and the cutoff, not $\gamma$.

\item \textbf{Three papers give an absolute scale, by different routes.}
Against an all-atom \emph{harmonic} model (Ming \& Wall): mechanical stiffness, independent
of temperature.
Against \emph{MD} variances or frequencies (Orellana; Dubanevics \& McLeish): effective
constants at the MD temperature, absorbing anharmonicity and finite sampling.

\item \textbf{For a uniform ANM, $\gamma$ is not a single number.}
It depends on the target and is coupled to the cutoff: 53.4 at 7.75~\AA\ (distribution
shape), 12--31 at 8--8.5~\AA\ (first MD mode), 0.4 at 15.4~\AA\ (MSDs), all in
kcal/(mol\,\AA$^2$).
Across proteins at fixed protocol it varies by a factor of about 2.5 (Dubanevics \&
McLeish).
The conventional $\gamma = 1$~kcal/(mol\,\AA$^2$) is none of these.

\item \textbf{Conventions differ by factors of 2.}
$\tfrac12 k\,\Delta d^2$ (Dubanevics \& McLeish; penm) versus $k\,\Delta d^2$ (Leioatts;
Orellana).
A constant taken from one paper into another code must be converted.

\item \textbf{B-factors are a poor calibration target.}
Ming \& Wall, Orellana and Fuglebakk all note that crystallographic B-factors include
static disorder, lattice effects and crystal constraints; Fuglebakk finds that fitting to
them selects networks with little collective motion.

\item \textbf{Shape is better settled than scale.}
Leioatts and Fuglebakk agree that what matters is a strong contrast between the first
neighbour shell (or sequence neighbours) and weaker, fast-decaying long-range springs; the
functional form matters little.
\end{enumerate}

\section{penm's models}

As of penm 0.6.0, every model returns spring constants in kcal/(mol\,\AA$^2$), so all
energies are in kcal/mol and \texttt{beta\_boltzmann()} (mol/kcal, CODATA gas constant)
applies to every model. Models published in kJ keep their published numbers in the code and
are converted on return.

\begin{center}
\small
\begin{tabular}{llp{0.42\textwidth}}
\toprule
\texttt{model} & published units of $k$ & source; in penm \\
\midrule
\texttt{ming\_wall} & kcal/(mol\,\AA$^2$) & \cite{ming2005}; default $k = 4.26$ (was 4.5 before 0.6.0) \\
\texttt{reach} & kJ/(mol\,\AA$^2$), MD at 300~K & Moritsugu et al.\ 2009, Table 1, solvated dimer; converted \\
\texttt{hnm} & kJ/(mol\,\AA$^2$) & Hinsen et al.\ 2000, as in bio3d \texttt{ff.calpha}; converted; C$_\alpha$ nodes only \\
\texttt{anm}, \texttt{gnm}, \texttt{pfanm}, \texttt{hnm0} & none & read as kcal/(mol\,\AA$^2$), arbitrary scale ($k = 1$, $d_{\max} = 10$~\AA\ for \texttt{anm}) \\
\bottomrule
\end{tabular}
\end{center}

For the arbitrary-scale models, energies and anything that adds an energy to an entropic
term carry that arbitrary scale; the calibrations of Section~2 are the alternatives.

\begin{thebibliography}{9}
\bibitem{ming2005} D.~Ming and M.~E.~Wall, Allostery in a coarse-grained model of protein
dynamics, \emph{Phys.\ Rev.\ Lett.} \textbf{95}, 198103 (2005).
\bibitem{orellana2010} L.~Orellana, M.~Rueda, C.~Ferrer-Costa, J.~R.~Lopez-Blanco,
P.~Chac\'on and M.~Orozco, Approaching elastic network models to molecular dynamics
flexibility, \emph{J.\ Chem.\ Theory Comput.} \textbf{6}, 2910 (2010).
\bibitem{leioatts2012} N.~Leioatts, T.~D.~Romo and A.~Grossfield, Elastic network models
are robust to variations in formalism, \emph{J.\ Chem.\ Theory Comput.} \textbf{8}, 2424
(2012).
\bibitem{fuglebakk2013} E.~Fuglebakk, N.~Reuter and K.~Hinsen, Evaluation of protein
elastic network models based on an analysis of collective motions, \emph{J.\ Chem.\ Theory
Comput.} \textbf{9}, 5618 (2013).
\bibitem{dubanevics2022} I.~Dubanevics and T.~C.~B.~McLeish, Optimising elastic network
models for protein dynamics and allostery: spatial and modal cut-offs and backbone
stiffness, \emph{J.\ Mol.\ Biol.} \textbf{434}, 167696 (2022).
\end{thebibliography}

\end{document}
